problem:  
Let \((M^m,g_M)\) be a complete, noncompact Riemannian manifold with Ricci curvature \(\mathrm{Ric}_M \ge 0\) and Euclidean volume growth  
\[
\mathrm{Vol}(B_R(x_0)) \simeq R^m,\quad R\to\infty.
\]  
Assume further that \(M\) supports a *sharp \(L^2\)–Sobolev inequality* of Euclidean type:
\[
\|v\|_{L^{2m/(m-2)}}^2 \le S_M\|\nabla v\|_{L^2}^2,\qquad v\in C_c^\infty(M).
\]

Let \((N^n,g_N)\) be a simply connected Cartan–Hadamard manifold with sectional curvature satisfying the two–sided bound
\[
-\,C_2 (1+r(x))^{-\alpha} \leq K_N(x) \leq -\,C_1 (1+r(x))^{-\alpha}, 
\qquad r(x)=d_N(p,x),
\]
for \(C_1,C_2>0\) and some decay exponent \(\alpha>0\).

Suppose \(u:M\to N\) is a harmonic map with finite energy
\[
E(u)=\frac{1}{2}\int_M |\nabla u|^2 < \infty.
\]

---

**Precise conjecture:**  
Let \(\alpha_c(m)\) denote the *critical curvature decay rate* separating rigidity from flexibility for such harmonic maps. Then:

- If \(\alpha < \alpha_c(m)\), every finite-energy harmonic map \(u:M\to N\) is constant.  
- If \(\alpha > \alpha_c(m)\), there exist examples of nonconstant finite-energy harmonic maps for some such \((M,N)\).

Determine whether the conjectural threshold
\[
\boxed{\alpha_c(m) = 2}
\]
is universal, i.e. whether quadratic curvature decay (\(|K_N|\sim (1+r)^{-2}\)) constitutes the boundary between global rigidity and potential nontriviality in the Euclidean-Sobolev scaling regime.

Equivalently: does negative curvature decaying slightly slower than \(r^{-2}\) still enforce constancy of all finite-energy harmonic maps from nonnegatively curved, Euclidean-growth manifolds?

---

Why is it a "good" problem:  
This formulation removes speculative functional dependence and focuses on a single, sharply defined analytic-geometric borderline: *quadratic curvature decay* on the target against *Euclidean Sobolev scaling* on the source. The question is simple to state—asking whether quadratic decay is the precise rigidity threshold—but profoundly deep: it probes the balance between curvature terms in the Bochner identity and the Sobolev-critical structure of energy decay.  

A resolution would reveal the exact asymptotic mechanism by which global geometry of \(N\) controls harmonic-map rigidity, uniting multiple classical vanishing results (constant-negative and subquadratic curvature) under a single critical principle. It also links curvature decay in differential geometry with analytic decay in elliptic PDEs, embodying the "narrow entrance, profound depth" ideal: a concise question that bridges Riemannian geometry, partial differential equations, and potential theory at the frontier of harmonic map research.